\documentclass[letterpaper,oneside,twocolumn,openany]{dndbook}

% Packages
\usepackage[utf8]{inputenc}
\usepackage[english]{babel}
\usepackage{lipsum}
\usepackage{listings}

% Document info
\title{The Shadow of Mars}
\author{A Roman Campaign in Five Acts}
\date{\today}

\begin{document}

\maketitle
\tableofcontents

\chapter{Introduction}

\section{Campaign Overview}

\textit{The Shadow of Mars} is a 5-session campaign set in a fantastical Roman world where the gods are real, auguries matter, and destiny can be changed—or embraced.

The party plays as members of a military unit stationed at a remote frontier fort. When dark omens appear and a powerful artifact is unearthed, they must decide whether to serve Rome, the gods, or their own ambitions.

\textbf{Levels:} Characters start at level 3 and advance to level 7 by the campaign's end.

\textbf{Themes:} Duty vs. ambition, fate vs. free will, loyalty vs. survival, divine intervention

\textbf{Setting:} The campaign takes place along Rome's Germanic frontier in 175 AD, in an alternate world where D\&D magic exists alongside Roman military discipline and pantheon worship.

\section{Key NPCs}

\begin{itemize}
\item \textbf{Legate Marcus Aurelius Corvinus} – Commander of the fort, ambitious and ruthless
\item \textbf{Augur Cassia Liviana} – The fort's seer, troubled by dark visions
\item \textbf{Centurion Titus Varro} – Veteran soldier, loyal but pragmatic
\item \textbf{The Tribune} – A mysterious figure from Rome with a hidden agenda
\item \textbf{Vercingetorix the Red} – Germanic chieftain and reluctant ally
\end{itemize}

\section{The Central Mystery}

A cursed artifact called the \textit{Spear of Mars} has been unearthed near the fort. Legend says it once belonged to the war god himself and was used to kill a divine being. Now it whispers to those who touch it, promising power and glory while slowly corrupting their souls.

The party must decide: destroy it, claim it, or deliver it to Rome?

\chapter{Session 0 Questions}

\section{Campaign Basics}
\begin{itemize}
\item Are you okay with this being a short campaign (about 5 sessions) with a clear ending?
\item Are character death and permanent consequences on the table?
\item Do you prefer more combat, more roleplay, or a balance of both?
\item Are you okay with the world reacting strongly to your choices?
\end{itemize}

\section{Tone \& Comfort}
\begin{itemize}
\item What tone do you want: serious, gritty, heroic, or light?
\item Are there any themes or topics you want to avoid?
\item Is PvP (stealing from or attacking party members) allowed?
\item Are you okay with NPC authority figures giving orders?
\end{itemize}

\section{Characters \& Party}
\begin{itemize}
\item Do your characters already know each other, or do they meet in session 1?
\item Why does your character choose to work with this group?
\item What is one clear strength your character brings to the party?
\item What is one weakness or flaw your character has?
\end{itemize}

\section{Combat \& Play Style}
\begin{itemize}
\item Do you prefer combat to be fast and cinematic or slower and tactical?
\item Are you okay with enemies retreating or surrendering?
\item Is retreat a valid option for the party?
\item Are you okay with teamwork and positioning mattering in combat?
\end{itemize}

\chapter{Chapter 1: Blood and Omens}

\section{Session Overview}
The party arrives at Fort Vindolanda on the Germanic frontier. Strange omens plague the camp, and the party is sent to investigate a nearby excavation where something ancient has been unearthed.

\textbf{Level:} 3 → Advance to 4 at session end

\textbf{Key Goal:} Investigate the excavation site and recover the artifact

\section{Opening Scene: The Fort at Dawn}

The party arrives at Fort Vindolanda as the sun rises over the Rhine. The wooden palisade walls are manned by weary legionaries, and smoke rises from the principia (headquarters). The camp is tense—yesterday, all the camp dogs died at once, and the augurs found no livers in the sacrificed sheep.

\begin{paperbox}{Read Aloud}
The frontier fort sprawls before you like a scar on the earth. Timber walls rise thirty feet high, crowned with watchtowers where sentries scan the dark forests beyond. The smell of wood smoke, leather, and sweat hangs in the air. As you pass through the gate, a legionary whispers: "The gods are angry. Turn back while you can."
\end{paperbox}

\section{Scene 1: Meeting the Legate}

The party is summoned to meet \textbf{Legate Marcus Aurelius Corvinus} in the principia. He's a hard man in his forties, scarred and ambitious.

\subsection{Legate's Briefing:}
\begin{itemize}
\item Three days ago, a work crew digging a new well discovered ancient ruins beneath the fort
\item They unearthed an iron spear covered in strange runes
\item That night, all the camp dogs died
\item The next day, the augurs found corrupted omens: black livers, reversed bird flights
\item One worker touched the spear and went mad, killing two companions before being subdued
\item The Legate wants the party to go down into the ruins, secure the artifact, and bring it to him
\end{itemize}

\subsection{Complications:}
\begin{itemize}
\item \textbf{Augur Cassia Liviana} interrupts, warning that the omens forbid disturbing the spear
\item \textbf{Centurion Varro} suggests destroying it immediately
\item The Legate overrules them both—Rome must possess this weapon
\end{itemize}

\textbf{Decision Point:} Will the party obey orders or heed the augur's warning?

\section{Scene 2: Descent into the Ruins}

The excavation site is a crude shaft dug through Roman stonework into older Germanic structures. The air grows cold as the party descends.

\subsection{Exploration:}
\begin{itemize}
\item Ancient Germanic temple dedicated to a forgotten war god
\item Walls covered in bronze age runes (DC 15 History or Religion to read)
\item Runes tell a story: \textit{"Here lies the God-Killer, the Spear of Broken Oaths, sealed by the blood of heroes. May it never taste sunlight again."}
\item Three chambers:
  \begin{enumerate}
  \item \textbf{Hall of Shields} – Dozens of ancient shields line the walls
  \item \textbf{Chamber of Chains} – Iron chains hang from ceiling, sized for something enormous
  \item \textbf{The Vault} – Central chamber where the spear was found
  \end{enumerate}
\end{itemize}

\subsection{Combat Encounter:} 2 Animated Armor (CR 1 each)
\begin{itemize}
\item The spirits scream: "DO NOT TAKE THE SPEAR!"
\item If defeated: They dissolve, whispering "You have doomed us all"
\end{itemize}

\section{The Spear of Mars}

\begin{monsterbox}{Spear of Mars}
\textit{Weapon (spear), legendary (requires attunement)}

The \textbf{Spear of Mars} rests on a stone altar, surrounded by broken chains. It's eight feet long, made of black iron, covered in glowing red runes.

\textbf{Mechanics:}
\begin{itemize}
\item Anyone who touches it must make a DC 14 Wisdom saving throw
\item \textbf{Failure:} Visions of glory and conquest flood their mind. They gain 1 level of corruption
\item \textbf{Success:} They resist but feel its pull
\end{itemize}

\textbf{Properties (if attuned):}
\begin{itemize}
\item +2 magic weapon
\item Deals an extra 1d8 necrotic damage
\item Disadvantage on Wisdom saving throws
\item User becomes increasingly aggressive and paranoid
\end{itemize}

\textbf{Whispers:} The spear whispers promises of heroism, conquest, and power.
\end{monsterbox}

\section{Corruption Mechanic}

Track corruption levels (0-5) for any character who touches or attunes to the spear:
\begin{itemize}
\item \textbf{Level 1:} Disturbing dreams of conquest
\item \textbf{Level 2:} Disadvantage on Wisdom saves
\item \textbf{Level 3:} Must succeed DC 12 Wisdom save to retreat from combat
\item \textbf{Level 4:} Must succeed DC 14 Wisdom save to avoid attacking allies when bloodied
\item \textbf{Level 5:} Fully corrupted, becomes NPC under DM control
\end{itemize}

\chapter{Chapter 2: The Tribune's Gambit}

\section{Session Overview}
A Tribune arrives from Rome with secret orders regarding the spear. Meanwhile, Germanic raiders attack the fort, and the party must defend it while uncovering the Tribune's true agenda.

\textbf{Level:} 4 → Advance to 5 at session end

\section{The Tribune Arrives}

Three days after the spear's discovery, a Tribune arrives from Rome with an elite escort of Praetorian Guards.

\begin{paperbox}{Read Aloud}
The gates open to admit a small column of riders in gleaming armor. At their head rides a man too clean for the frontier—purple-trimmed toga, oiled hair, rings on every finger. He surveys the fort like a merchant appraising damaged goods. Behind him, six Praetorians sit their horses like statues of iron and menace.
\end{paperbox}

\textbf{Tribune Lucius Valerius Maximus} introduces himself with imperial authority and relieves the Legate of command.

\section{The Raid}

That night, Germanic raiders attack the fort in force—coordinated with the Tribune's arrival (suspicious timing).

\subsection{Combat Phases:}

\textbf{Phase 1: The Walls (2 rounds)}
\begin{itemize}
\item Raiders swarm the palisade with ladders
\item Party can man ballistae (2d10 piercing, DC 12 Dex to aim)
\end{itemize}

\textbf{Phase 2: Breach (3 rounds)}
\begin{itemize}
\item Raiders break through the gate with a battering ram
\item Melee combat in the fort courtyard
\end{itemize}

\textbf{Phase 3: The Headquarters (2 rounds)}
\begin{itemize}
\item \textbf{Vercingetorix the Red} and berserkers attack
\item Climactic battle to defend the principia
\end{itemize}

\chapter{Chapter 3: Through the Dark Forest}

\section{Session Overview}
The party must journey through the Teutoburg Forest—either fleeing Rome or on a secret mission to destroy the spear.

\textbf{Level:} 5 → Advance to 6 at session end

\begin{paperbox}{Read Aloud}
The forest swallows you whole. Ancient oaks rise like cathedral columns, their canopy so thick that even at midday, twilight reigns. The Romans call this the Dark Forest—the place where legions disappear.
\end{paperbox}

\section{The Sacred Grove}

Deep in the forest lies an ancient Germanic holy site where the spear can be destroyed.

\subsection{The Ritual Requirements:}
\begin{enumerate}
\item The spear must be placed on the altar
\item Blood of a warrior must anoint it (1d6 damage, voluntary)
\item A sacrifice must be made—something precious given willingly
\item The party must speak the words: "We unmake what was made in darkness"
\end{enumerate}

\section{Mars Intervenes}

As the ritual begins, \textbf{Mars himself manifests}—not physically, but as a divine presence.

\begin{commentbox}{Encounter: Avatar of Mars}
Use \textbf{Fire Elemental} stat block (CR 5) but reskin as divine wrath. Mars speaks: "THE SPEAR IS MINE! YOU DARE DESTROY WHAT I CREATED?"
\end{commentbox}

\chapter{Chapter 4: Return to Rome}

\section{Session Overview}
The party returns to Rome as either heroes or fugitives. They must navigate the political intrigue of the Senate and uncover who truly wanted the spear.

\textbf{Level:} 6 → Advance to 7 at session end

\section{The Conspiracy Revealed}

\textbf{Senator Gaius Cassius Brutus} is revealed as the mastermind behind the plot to acquire the spear and assassinate the Emperor.

\subsection{The Underground Temple}

\begin{paperbox}{Read Aloud}
Beneath the city, deeper than the sewers, you find it—a cavern carved into living rock. The walls run red with ochre paint made to look like blood. An altar of black stone dominates the center, and upon it... a replica spear. Around it, cultists chant in Latin: "Mars Ultor! Mars Ultor! Avenge us!"
\end{paperbox}

\chapter{Chapter 5: The Wrath of Mars}

\section{Session Overview}
The final confrontation. Mars demands satisfaction for the destruction of his spear. The party must face divine judgment.

\textbf{Level:} 7 (final level)

\section{The Divine Summons}

\begin{paperbox}{Read Aloud}
The sky turns red at noon. Not sunset red—blood red. The Tiber runs backwards. Every dog in Rome howls in unison. In the Forum, the statue of Mars begins to weep crimson tears.

That night, you all share the same dream. You stand in a vast arena of black sand beneath a blood-red sky. Mars himself sits upon a throne of swords and bones. His voice shakes your souls: "YOU DESTROYED WHAT WAS MINE. NOW PAY THE PRICE."
\end{paperbox}

\section{Trial by Combat}

Mars offers three options:
\begin{enumerate}
\item One party member fights Mars in single combat
\item Entire party fights Mars together
\item Attempt to bargain with the god
\end{enumerate}

\section{The Price of Victory}

After defeating or bargaining with Mars, there is always a cost:

\begin{itemize}
\item \textbf{Marked by War:} Each character bears Mars' brand
\item \textbf{The Burden:} One becomes Mars' Chosen
\item \textbf{The Debt:} They owe Mars a future favor
\item \textbf{The Prophecy:} "In three generations, Rome will burn"
\end{itemize}

\chapter{Epilogue}

\section{Campaign Conclusion}

The party awakens in the Temple of Mars in Rome. Their choices throughout the campaign determine their ending:

\begin{itemize}
\item \textbf{Heroic Ending:} Hailed as saviors of Rome
\item \textbf{Pragmatic Ending:} Quietly rewarded, leave Rome
\item \textbf{Tragic Ending:} Some died in Mars' arena
\item \textbf{Rebellious Ending:} Hunted beyond the Empire
\end{itemize}

\begin{paperbox}{DM Final Narration}
The Shadow of Mars is long. What you did in those five sessions echoes across the Empire. Some call you heroes. Some call you fools. Mars himself calls you... interesting. And in the divine realm, that is far more dangerous than being forgotten. Your story is over. But legends never truly end.
\end{paperbox}

\section{Guy Sclanders' Principles Applied}

\subsection{Master Plot}
The entire campaign followed the "Spear's Journey": Discovery → Corruption → Destruction → Divine Judgment. Each session escalated the stakes while maintaining player agency.

\subsection{Theme}
\textbf{Duty vs. Free Will} - The party was caught between Rome's orders, the gods' will, and their own conscience.

\subsection{Villain Archetypes}
\begin{itemize}
\item \textbf{Legate Corvinus:} Blunt Force Nemesis (direct opposition, Session 1-2)
\item \textbf{Senator Brutus:} Never Present (manipulating from shadows, revealed Session 4)
\item \textbf{Mars:} Divine Antagonist (ultimate challenge, Session 5)
\end{itemize}

\chapter{Appendix: Quick Reference}

\section{Key NPCs Stats}

\begin{itemize}
\item \textbf{Legate Marcus Corvinus:} Use \textbf{Knight} (CR 3)
\item \textbf{Tribune Lucius:} Use \textbf{Noble} with +2 AC (CR 2)
\item \textbf{Centurion Varro:} Use \textbf{Veteran} (CR 3)
\item \textbf{Augur Cassia:} Use \textbf{Priest} (CR 2)
\item \textbf{Vercingetorix:} Use \textbf{Gladiator} (CR 5)
\item \textbf{Senator Brutus:} Use \textbf{Priest} + magic armor (CR 3)
\end{itemize}

\section{Roman Military Ranks}

\begin{itemize}
\item \textbf{Legate} – Commander of the legion (like a general)
\item \textbf{Tribune} – Staff officer from Rome (aristocrat)
\item \textbf{Centurion} – Company commander (60-80 men)
\item \textbf{Optio} – Second-in-command to centurion
\item \textbf{Legionary} – Standard soldier
\end{itemize}

\vfill
\begin{center}
\textbf{Campaign Complete! May Mars smile upon your future games!}

\vspace{1em}
\textit{Created with Guy Sclanders' methodologies}

\textit{Version 1.0 - January 2026}
\end{center}

\end{document}
